\section{Parámetros de comparación}\label{sec:parametros}

Las baterías se comparan con los parámetros que aparecen en sus dentro de sus hojas de datos y con base a sus especificaciones se dimenciona el banco. La Tabla~\ref{tab:parametros} resume qué indica cada uno y cómo se usa en el diseño.

La \emph{capacidad nominal} es la carga energetica que la batería entrega al descargarse a una corriente constante hasta un valor de  tensión final, en condiciones definidas por el fabricante \cite{linden}. Se expresa en Ah o, multiplicada por la tensión nominal, en Wh. Como la capacidad de la bateria disminuye cuando la corriente de descarga aumenta, el fabricante la declara para un tiempo de descarga: $C_{20}$ es la capacidad cuando la batería se descarga en \SI{20}{h}, y $C_{100}$ cuando se descarga en \SI{100}{h} \cite{linden}. En sistemas fotovoltaicos se usan normalmente $C_{20}$ y $C_{100}$, porque la descarga ocurre a lo largo de la noche con corrientes bajas. Para comparar modelos entre sí se usa la misma tasa en todos.

La \emph{profundidad de descarga} (DOD) es la fracción de la capacidad que se extrae en un ciclo, y el \emph{estado de carga} (SOC) es su complemento, $\mathrm{SOC} = 1 - \mathrm{DOD}$. La DOD se encarga de fijar qué parte de la capacidad nominal se puede usar y, a la vez, cuántos ciclos resiste la batería: entre más profunda es la descarga, menor es el número de ciclos \cite{linden}. Por eso los \emph{ciclos de vida} solo se pueden comparar si están declarados a la misma DOD.

La \emph{eficiencia} de ida y vuelta es la relación que se tiene entre la energía que la batería entrega y la que recibió para cargarse. La \emph{autodescarga} es la pérdida de carga de una batería que se encuentra en reposo, y se expresa como porcentaje por mes. La \emph{temperatura} corresponde a un parametro que afecta la capacidad disponible y la vida útil, y cada tecnología tiene un rango distinto para cargar y para descargar. El \emph{mantenimiento} indica si la batería depende de reposición de agua, cargas de ecualización o ventilación cada cierto tiempo. Además, la corriente máxima de carga y de descarga limita la potencia que el banco puede recibir del regulador y entregar al inversor, y puede obligar a conectar baterías en paralelo aunque la capacidad ya sea suficiente.

El \emph{costo} se compara por kWh útil, es decir, por la energía que se puede extraer en un ciclo sin superar la DOD recomendada:
\begin{equation}
c_{\text{útil}} = \frac{1000 \times \text{precio}}{V_{\mathrm{nom}}\,C\,\mathrm{DOD}},
\label{eq:costo_util}
\end{equation}
en USD/kWh, donde $V_{\mathrm{nom}}$ es la tensión nominal en V y $C$ la capacidad en Ah. Una batería barata por Ah puede resultar cara por kWh útil si su DOD recomendada es baja.

\begin{table}[!ht]
\caption{Parámetros de comparación de las baterías}
\label{tab:parametros}
\centering
\footnotesize
\setlength{\tabcolsep}{3pt}
\begin{tabular}{>{\raggedright\arraybackslash}p{0.27\columnwidth} >{\raggedright\arraybackslash}p{0.66\columnwidth}}
\toprule
Parámetro & Uso en el diseño \\
\midrule
Capacidad nominal (Ah, Wh) & Número de baterías en paralelo; se compara a la misma tasa de ($C_{20}$ o $C_{100}$). \\
Tensión nominal (V) & Número de baterías en serie para formar \SI{24}{V}. \\
DOD recomendada (\%) & Fracción utilizable de la capacidad; costo por kWh útil. \\
Ciclos de vida & Años hasta el reemplazo, a la DOD de diseño. \\
Eficiencia (\%) & Pérdidas en la batería dentro del balance de energía. \\
Autodescarga (\%/mes) & Pérdida en reposo; factor $K_a$ del balance. \\
Corriente máx. de carga y descarga (A) & Baterías en paralelo según el regulador y el inversor. \\
Temperatura de operación (°C) & Capacidad disponible y condiciones del lugar de instalación. \\
Mantenimiento & Agua, ecualización, ventilación, sistema de gestión (BMS). \\
Peso (kg) y garantía & Instalación en la casa y vida útil asegurada. \\
Costo (USD/kWh útil) & Comparación económica, ecuación~\eqref{eq:costo_util}. \\
\bottomrule
\end{tabular}
\end{table}
